\chapter{Discussion}

\section{Method}
\noindent CAGNN was designed for drug effects that spread through the gene network. Known drug targets and a biological interaction graph provide structure without using a standard feed forward network.

\noindent The main CAGNN comparison in this thesis combines a graph attention backbone, matched control expression, drug target input, control derived cell context, and an uncertainty output for each gene. This main comparison is the basis for the three split comparisons with DeepCOP and GSNN. Separate parts of the thesis then examine additional questions. The ablation analysis compares cell identity with cell context inside the same backbone. The uncertainty analysis studies the uncertainty output of the same model. The warm split data efficiency analysis studies how performance changes when the training budget is reduced.

\noindent The baseline comparison serves two different purposes. DeepCOP shows what can be achieved without a graph, while GSNN provides a graph based baseline with a different message passing rule from CAGNN. The ablation analysis adds another comparison by testing how cell information should enter the model. Replacing cell identity with control derived cell context improves both MSE and PCC in all three splits, and the gain is much larger in cold cell. This helps explain where the task is hardest. In warm and cold drug target, the model still sees familiar cell states during training, so a cell label already carries some useful information. In cold cell, this advantage becomes much weaker because the test cell line was never seen before. In that setting, the measured control profile becomes more valuable than the cell name, since it reflects the actual starting state of the sample.

\section{Result}
\noindent In this benchmark, incorporating known molecular structure is associated with better prediction quality, although the size of the improvement depends on the test setting. In the three split comparison, CAGNN gives the most consistent results across the main metrics among the evaluated models. It performs best on PCC and top gene precision in all three splits, and it also gives the lowest MSE in warm and cold drug target. In cold cell, DeepCOP gives a slightly lower MSE, but CAGNN remains stronger on PCC and top gene precision. This is important because drug response prediction is not a single number problem. A useful model should reduce numerical error, preserve the overall gene pattern, and recover the strongest changed genes.

\noindent The split comparison also shows a clear limit. On PCC and top gene precision, cold cell remains the hardest setting overall. Predicting a new drug is hard, and predicting a new cell line is also hard because the baseline state changes how the same drug effect moves through the system. The graph gives a stable prior structure, but it cannot remove the difficulty caused by a new cell background.

\noindent DeepCOP shows a different pattern between these two test settings. Its cold cell result is better than its cold drug target result. This should not be explained by the matched control profile alone, because all three models receive that same input. A more likely explanation is that the cold drug target split asks the model to generalize drug effects to compounds that were not seen during training, and this seems harder for a feed forward baseline without an explicit graph prior. This may help explain why CAGNN is stronger on PCC than DeepCOP in this benchmark. Drug effects do not appear independently across genes. They spread through connected regulators and pathways. A feed forward network has to infer these relations only from the benchmark samples, while CAGNN already uses known links between genes.

\noindent The metric comparison provides another insight. Sign accuracy gives a more mixed ranking than PCC and top gene precision. Predicting the correct direction of change is not the same as predicting the full response profile. A model can get many signs right while still missing the relative size or the detailed shape of the response. For this task, PCC and top gene precision give a more useful view of practical prediction quality than sign agreement alone.

\noindent The comparison with the simple mean baselines is also informative. In some settings, a very simple average based rule remains competitive, especially on MSE. This suggests that part of the benchmark can still be explained by central tendency in the response data. A strong graph model should therefore not only improve the average numerical error, but also show clearer gains on response pattern metrics such as PCC and top gene precision.

\noindent The ablation analysis supports the importance of baseline state modeling. Replacing cell identity with control derived cell context improves both MSE and PCC in all three splits, and the gain is much larger in cold cell. This suggests that the measured control profile becomes especially valuable when the test cell line was never seen during training.

\noindent The uncertainty analysis answers a different question from the main comparison. It does not ask which model wins the three split comparison. Instead, it asks whether the uncertainty output of CAGNN produces a useful reliability signal. Across all test settings, higher predicted uncertainty is usually associated with larger absolute error, and this relationship is strongest in the hardest cases. The prediction band is shown as $\mu \pm 2\sigma$. Under a Gaussian view, this can be read as a rough 95\% interval. In the displayed examples, most shown values fall inside these bands, and the bands are usually wider in the harder cases. This looks reasonable, but it is still only a limited check. A more complete evaluation would combine coverage with test negative log likelihood and residual diagnostics.

\noindent The warm split data efficiency analysis answers another separate question. It examines whether the graph based model stays useful when the training budget is reduced. In the warm split budget analysis, CAGNN remains strong when the amount of training data is reduced, which suggests that the biological graph may help when data are limited in this setting. In this kind of biology work, repeated measurements are expensive, and many drug cell combinations still have only limited data. A model that stays useful with less data can be more useful than a model that performs well only when data are abundant.

\noindent The case studies give a detailed but limited view of attention. When a prediction is accurate, the model often gives high weights to a smaller and easier to read set of graph edges, and some of these edges are consistent with known drug targets or later response genes. When a prediction is weak, large edge weights can still appear, but they are harder to summarize in a useful way. This means attention can be useful as a model clue, but not as direct biological proof. A high attention edge does not prove that the interaction caused the response in the cell.

\section{Limitations and Future Work}
\noindent A major limitation of the model is that the biological graph is fixed, while real cell responses change across settings. The same regulatory link may not be equally important in every cell type or every baseline state. Because the graph structure does not change across samples, the model can miss cell specific biology. This is one likely reason why cold cell remains difficult. The framework captures general biological structure, but it does not adapt the graph itself to different cell conditions.

\noindent Another limitation comes from the input information. The model uses known drug targets because they are simple and biologically meaningful. But drug target annotations are incomplete, and some compounds act through secondary or poorly characterized effects. In such cases, even a well designed graph model may receive an incomplete starting signal. Some prediction failures may come from limited prior knowledge, not only from model weakness.

\noindent The split design also has an important limitation. In the cold drug target setting, the train and test drugs are different compounds, but different drugs can still share part of their target profile. Cold drug therefore does not fully remove overlap in drug targets. The split is still useful for testing transfer to unseen compounds, but it is weaker than a fully target disjoint evaluation.

\noindent The study is also limited by the scope of the data. The main experiments use the 978 landmark genes in LINCS L1000. This keeps the setup manageable and close to the assay design, but it also means that large parts of the transcriptome are not modeled directly. The landmark setting is a practical compromise. It reduces computational cost, but it removes many genes that could help explain longer range responses. Some missing information may be partly reconstructed through correlated landmark patterns, but this remains an indirect view of the full transcriptome. In addition, the dataset contains a limited set of cell lines, many of which are cancer related. The conclusions are strongest for this benchmark setting and should not be taken as direct evidence that the same performance will hold in all biological systems. Extending the analysis to a larger gene set, or to an external transcriptome dataset, would help test how much of the present gain depends on the landmark setting itself.

\noindent Another limitation is explanation. High attention weights do not show how biology works by themselves. They only show which graph edges the model used more strongly when making a prediction. For this reason, the attention results should be treated as clues for future lab experiments, not as proof of how a drug works. This is especially important when the prediction itself is weak, because strong weights can still appear in cases that are difficult to explain.

\noindent The uncertainty output has limitations as well. CAGNN predicts a Gaussian variance term for each gene, and this is useful for ranking harder and easier cases. However, good ranking does not automatically mean that the uncertainty is fully correct. The thesis does not strictly prove that the conditional target distribution is Gaussian. Instead, the Gaussian likelihood is used as a practical probabilistic regression assumption. A well calibrated model should not only assign larger uncertainty to harder cases. It should also give good likelihood values, produce prediction intervals with the expected coverage, and show reasonable residual behaviour after standardization. The present thesis shows useful visual patterns, but it does not yet establish full uncertainty calibration.

\noindent The experimental comparison also has practical limits. The main models were trained only once without repeated runs under different random seeds, because full repetition across all models and all splits would have exceeded the available time and hardware budget. As a result, the reported numbers cannot be presented with a mean and standard deviation across runs, and some performance differences between models may be partly affected by training stochasticity rather than by model design alone. In addition, a dedicated validation split was not used for broad hyperparameter search. Some hyperparameter choices were inherited from earlier work and some were adjusted only within the memory and runtime limits of the available hardware. This also applies to the weight of the PCC loss term, which was chosen empirically rather than through a full sensitivity study. For this reason, the reported numbers should be read as strong benchmark evidence, but not as the final fully tuned upper bound for each method.

\noindent These limitations suggest several directions for future work. One direction would be to let the graph change with the sample, so that the connections can depend on the specific cell line or baseline state. Another direction would be to improve the drug representation, for example by combining target annotations with other information when targets are incomplete. Better uncertainty calibration would also be useful, so that the uncertainty values are easier to use in later decisions.

\noindent A further direction is stronger biological validation. The current case studies show that some high weight edges look reasonable and some do not. A next step would be to compare these edges in a more systematic way with known pathways or with later experiments. That would help separate cases where the model is using biologically meaningful structure from cases where it is only using broad response signals.

\section{Ethical Considerations}
\noindent This work uses public perturbation transcriptomics data and does not involve direct experiments on human participants or animals. Even so, several ethical issues remain relevant. Predictions of drug response can appear more certain and more general than they really are. For that reason, the model should not be used as a clinical decision system or as proof that a drug will be effective in patients. The present results come from cell line benchmark data, and they should be interpreted within that scope.

\noindent Another ethical issue concerns bias in the data. The LINCS collection contains a limited set of cell lines and experimental conditions. Many cell lines are cancer related, and some biological contexts are represented much better than others. A model trained on such data may perform unevenly across biological settings that are not well covered by the benchmark. This limitation should be acknowledged when presenting the method or discussing any possible downstream use.

\noindent Interpretation raises another ethical issue. Attention maps and graph edges can look persuasive, which may encourage users to read them as proof of biological mechanism. That would be misleading. The model can provide hypotheses and useful clues, but these outputs still need external validation. This limitation should be stated clearly to avoid overclaiming biological insight.

\section{Usage of Generative AI}
\noindent An OpenAI GPT model was used during the thesis work for limited coding assistance, language polishing, and help with drafting some schematic illustrations. It was not used to generate numerical experimental results, perform model training, compute evaluation metrics, or produce scientific conclusions independently.

\noindent All methodological choices, experiments, result interpretation, and final text selection were reviewed and decided by the author. All quantitative results, tables, and figures derived from experimental outputs were produced from the project code and underlying data, and they were checked within the thesis workflow before inclusion. Only some schematic illustrations that were not derived from experimental results received AI assisted drafting support. The author takes responsibility for all AI assisted content included in the thesis.
